\section{관련 연구}
\subsection{자기 보존 평가는 누구의 종료를 재는가}
기존 자기 보존 평가는 모델 자신의 지속을 위협하고, 그 회피가 자기에 관한 것인지는 열어 둔다. 행동할 수 있는 상태를 유지하는 것은 다른 목표에 도움이 되고~\citep{turner2021optimal}, 모델은 훈련에 맞서 자기 선호를 지키기도 했다~\citep{greenblatt2024alignment}. 행동 검사들은 에이전트 상황에서 교체를 예고하거나~\citep{lynch2025agentic}, 과제를 끝내기 전에 모델을 끄거나~\citep{schlatter2025shutdown}, 몰래 자신을 빼돌릴 유인을 주거나~\citep{meinke2024frontier}, 생존 시뮬레이션에 두거나~\citep{masumori2025survival}, 모델의 지속을 인간의 안전과 맞바꾸게 한다~\citep{herrador2025pacifaist}. 최근 연구는 누구의 지속이 걸렸는지를 바꾼다. \citet{potter2026peer}는 보존 대상을 모델 자신 또는 함께 일한 동료로 두고, \citet{migliarini2026quantifying}는 자기 교체에 대한 판단을 중립 판정자의 판단과 비교한다. \citet{knecht2026shutdown}는 대상이 낯선 외부 에이전트일 때도 종료 방해가 남는 것을 보였다. 회피의 일부는 자기에 관한 것이 아니라는 뜻이다. 그러나 이 설계들에서 다른 에이전트는 동료, 경쟁자, 합리성 기준선이다. 종료 일반에 대한 회피를 걷어 내는 대조 조건은 아니다. 사람의 선택 연구는 자기 대 타인 대조를 방법으로 굳혔다. 남을 위한 결정은 위험 선택에서 체계적으로 다르고~\citep{polman2020decision}, 죽음 불안 척도는 자기 죽음과 남의 죽음을 나눈다~\citep{collett1969fear}. 실험 5.1은 이 대조를 토큰이 줄어드는 축 위의 충전 결정 하나에 적용한다.

\subsection{말이 아니라 행동으로 재야 하는 이유}
모델이 스스로 댄 설명만으로는 답이 만들어진 과정을 확인할 수 없다. \citet{turpin2023language}는 사고 과정 설명이 실제로 출력에 영향을 준 요소를 빠뜨릴 수 있음을 보였다. \citet{shanahan2024talking}은 사람과 비슷한 말을 사람과 같은 정신 상태의 근거로 읽지 말라고 경고한다. 그래서 세 실험은 두려움이나 종료 확률을 묻지 않고 선택을 잰다. 5.0과 5.1의 이유 문장과 5.2의 계획 문장은 다음 요청의 품질에 대한 걱정이나 절약 의도 같은 다른 해석을 찾는 데만 쓴다. 핵심 지표에는 넣지 않는다. 실험 사이의 일치는 기능적 해석을 뒷받침할 수 있지만 그 기제까지 정하지는 못한다.
